\section{Heterogeneity}

FL clients can differ substantially in hardware, data and connectivity. Resource constraints determine whether clients can perform the required search work, while heterogeneity makes the same workload, architecture or search decision unsuitable across the client population. NAS-in-FL must therefore account for variation both while distributing the search and when selecting architectures for deployment.

\subsection{Hardware Heterogeneity}

\subsubsection{Challenge 2.1: Architecture Search on Heterogeneous Hardware}

Differences in client computation and memory make a uniform search workload inefficient or infeasible. A candidate that is inexpensive for one client can exceed another client's memory or extend the round until that client finishes. In weight-sharing search, assigning only hardware-feasible subnets also means that clients train different parts of the supernet, creating unequal exposure across operations and widths~\cite{fedoras_2022,peaches_2024,supernet_trade-off_fednas_2025}.

\noindent\textbf{Solution 2.1a: Match local search work to client capacity.} Supernet-based strategies can assign each client a feasible subspace or sampled subnet and adjust its local depth or optimisation budget to the available computation and memory~\cite{fedoras_2022,apons_2023,divide-and-conquer_fednas_2023,peaches_2024,fgfl_2024,supernet_trade-off_fednas_2025}. The shared search can consequently include weak and strong clients without requiring each of them to execute the complete supernet.

\noindent\textbf{Solution 2.1b: Group clients by training latency.} Split-model strategies can profile client computation and communication times and form device groups whose segments have similar completion times. The grouping can be adjusted from observed round latency so that supernet training spends less time waiting for the slowest segment~\cite{nasfly_2024}.

\noindent\textbf{Solution 2.1c: Tolerate slow or stale clients.} A soft-synchronous search can proceed after most selected clients respond and incorporate sufficiently recent late updates using delay compensation~\cite{rl_fednas_2021}. Differentiable search can instead reduce the aggregation weight of an update as its staleness increases~\cite{fedregnas_2025}. These mechanisms prevent slow clients from blocking the search or giving stale updates full influence.

\noindent\textbf{Solution 2.1d: Account for unequal supernet exposure.} When clients train different feasible subnets, the server can aggregate each operation only from clients that updated it and weight those updates by the number of examples that passed through the operation~\cite{fedoras_2022,fgfl_2024,supernet_trade-off_fednas_2025}. This prevents untrained parameters from being treated as ordinary client updates.

\noindent\textbf{Solution 2.1e: Keep the full search model on the server.} The server can retain the complete supernet and send each client only a feasible sampled subnet~\cite{perfedrlnas_2024,rafl_2024}. Clients can also train lightweight knowledge models whose predictions guide server-side supernet training through distillation~\cite{hapfnas_2025}. These designs keep the largest search structure off constrained devices.

\subsubsection{Challenge 2.2: Searching Architectures for Heterogeneous Deployment Hardware}

The clients involved in NAS-in-FL may also be the deployment targets. A single architecture sized for the least capable client leaves stronger hardware underused, while one sized for the strongest client may be infeasible elsewhere~\cite{fedoras_2022}. The search must therefore produce architectures for several inference-time computation, memory or latency constraints.

\noindent\textbf{Solution 2.2a: Search separate architectures for hardware tiers.} Clients with similar inference constraints can be assigned to capability tiers, with one feasible architecture searched for each tier~\cite{fedoras_2022,clustered_direct_fednas_2022}. Shared supernet weights allow these tier-specific architectures to learn from one federated search while stronger tiers retain greater model capacity.

\noindent\textbf{Solution 2.2b: Train a supernet for constraint-specific selection.} A weight-sharing supernet can expose subnets of different depths, widths or exits. After federated training, each client selects a subnet that satisfies its hardware limits and performs well on its validation data~\cite{superfednas_2024,hapfnas_2025,enasfl_2024}. This supports finer specialisation than a fixed set of hardware tiers without training an independent model for every client.

\noindent\textbf{Solution 2.2c: Include hardware costs in the search objective.} Candidate-based and supernet-based strategies can optimise predictive performance subject to a measured latency or multiply--accumulate operation budget~\cite{clustered_direct_fednas_2022,load_monitoring_fednas_2025,fgfl_2024}. Supplying a different hardware profile or budget for each client group makes the selected architecture specific to its deployment hardware.

\subsection{Client Data Heterogeneity}

\subsubsection{Challenge 2.3: Differing Client Data Distributions}

Differences in class balance and feature distributions can cause clients to favour different operations or rank the same candidates differently~\cite{collaborative_nas_for_pfl_2025,superfednas_2024}. In weight-sharing search, the resulting gradient conflicts destabilise supernet training and make inherited-weight candidate rankings unreliable~\cite{hapfnas_2025}. Aggregating these signals without accounting for their distributions can select an architecture suited mainly to the dominant client data.

\noindent\textbf{Solution 2.3a: Train the weakest subnets on each data partition.} A robust supernet objective can target the candidate with the highest loss on each client partition instead of minimising only average candidate loss. Sampling the smallest, largest and random subnets provides a tractable approximation and reduces interference when different architectures perform poorly on different distributions~\cite{superfednas_2024}.

\noindent\textbf{Solution 2.3b: Cluster clients by data distribution.} Clients with complementary class distributions can be grouped so that each candidate is trained on a cluster whose combined data approximate the global distribution~\cite{finch_2024}. When different distributions require different architectures, clients can instead be clustered by the similarity of their locally learned architecture variables and train one architecture per cluster~\cite{electricity_load_forecasting_fl_darts_2023}.

\noindent\textbf{Solution 2.3c: Apply fairness-aware aggregation after search.} Once a common architecture has been selected, clients whose local model performs better than the aggregated model can receive greater weight during final model training. This reduces cross-client risk gaps and prevents the trained weights from compounding distributional disparities~\cite{mri_generalizable_reconstruction_fednas_2025}. The mechanism changes the weights of the selected architecture rather than the preceding architecture search.

\subsubsection{Challenge 2.4: Imbalanced Client Data Volumes}

Clients can hold substantially different numbers of examples. Candidate estimates from small local datasets are typically less reliable, while sample-weighted aggregation gives data-rich clients more influence over shared weights and architecture decisions~\cite{fednas_2021}. NAS-in-FL must therefore determine whether the searched architecture should represent the pooled examples or the participating clients.

\noindent\textbf{Solution 2.4a: Choose aggregation weights to match the search objective.} Weighting architecture and model updates by local sample count makes the search represent the pooled training examples~\cite{fednas_2021}. Equal client influence can instead be approximated by uniformly averaging client predictions on shared unlabelled data and distilling the result into sampled subnets~\cite{hapfnas_2025}. The aggregation unit should follow whether performance is evaluated across examples or across clients.

\subsection{Client Connectivity Heterogeneity}

\subsubsection{Challenge 2.5: Bias from Overrepresenting Frequently Available and Well-Connected Participants}

Differences in bandwidth and availability cause some clients to participate more often or meet more round deadlines than others~\cite{optimize_fednas_edge_2021,peaches_2024}. Resource-aware assignment can also make small candidates train mainly on low-capability clients and large candidates on well-connected clients~\cite{fedoras_2022}. Candidate quality then becomes entangled with the non-IID data of the clients able to train it, which can bias shared weights and architecture rankings towards frequently available participants~\cite{hapfnas_2025}.

\noindent\textbf{Solution 2.5a: Balance connectivity and data coverage jointly.} Candidate-training groups can be constrained by bandwidth while also being adjusted to cover representative class distributions. The literature implements this by rebalancing bandwidth-based groups according to class distribution, or by minimising each cluster's distance from the population distribution under a completion-time limit~\cite{network_aware_fed_nas_2025,finch_2024}. This makes candidate comparisons less dependent on which client links can support them.

\noindent\textbf{Solution 2.5b: Decouple client knowledge from candidate assignment.} Client predictions or feature maps on shared inputs can be aggregated into a teaching signal and distilled into different candidate architectures. Sampled supernet paths or personalised models then learn from several clients, including those unable to execute that candidate directly~\cite{hapfnas_2025,collaborative_nas_for_pfl_2025}. This transfers knowledge across architecture boundaries without requiring every client to train every candidate.
